% CVPR review-mode manuscript. All numbers are generated from
% results/p3_p4/metrics_completed.json; see paper/tables.tex and figures/.
\documentclass[10pt,twocolumn,letterpaper]{article}
\usepackage[review]{cvpr}
\input{preamble}
\usepackage{booktabs}
\usepackage{amsmath,amssymb}
\usepackage{graphicx}
\definecolor{cvprblue}{rgb}{0.21,0.49,0.74}
\usepackage[pagebackref,breaklinks,colorlinks,allcolors=cvprblue]{hyperref}

\def\paperID{*****}
\def\confName{CVPR}
\def\confYear{2027}

\title{Path-Aware Controls Decide Structure Effects in\\ Connectome-Constrained Models}

\author{Anonymous}

\setlength{\floatsep}{3pt plus 1pt minus 1pt}
\setlength{\textfloatsep}{4pt plus 1pt minus 1pt}
\setlength{\intextsep}{3pt plus 1pt minus 1pt}
\begin{document}
\maketitle

\begin{abstract}
Reports disagree on whether connectome topology benefits visual computation,
but they use different null families, horizons and tasks and are not directly
comparable. We preregister a nested null ladder with machine-checked invariants
and a path-identifiable two-hop XOR task, and show that matching degree, module
and triad statistics cannot localise an effect carried by a convergence motif;
we then carry the diagnosis onto real MaleCNS v1.0 wiring. \emph{Scope}: our
claims hold for controlled low-in-degree graphs and a specific real MaleCNS
subgraph regime, where path-aware control construction decides the observed
structure effect and the published weak/strict ladder is reproduced on that
real data. The preregistered \texttt{mechanism\_transfer\_confirmed=false}
means we cannot claim a causal mechanism for the literature disagreement, and
topology specificity is negative on real visual tasks.

We report distinct preregistered results, neither merged nor quoted
selectively. First, the fixed-density path-survival dose: the original C1
failed (\texttt{hypothesis\_supported=false}) because degree-compensating
swaps created fresh two-hop paths despite every global invariant passing; C1b
is an \emph{adaptive, remedial} preregistration, frozen before training but
designed after inspecting C1's in-flight readings, and its path-blocking
repair yields a monotone dose satisfying all six rules
(\texttt{hypothesis\_supported\_pathblock=true}) -- a new test, not
confirmation of the original hypothesis. Second, a faithful rerun of the
published control ladder reproduces the apparent early advantage and its
strict abolition but not its task-aware recovery, so
\texttt{mechanism\_transfer\_confirmed=false}; a connectome-constrained model
gains from having a graph, not the specific topology, and a clean
hash-verifiable main-matrix rerun passes all six gates. The conjunctive
\texttt{expected\_strong} target is \emph{not} met overall. The independent
24-graph family (C4) is now complete 24/24 and is the first real test of
contract-v2's second falsifiable prediction; it returns a mixed, partly
negative verdict \texttt{independent\_family\_validated=false}. All 24 deltas
are non-negative (eight high-$k$ graphs $\approx+0.50$, the rest $\approx0$),
so LOGO sign agreement $0.5000$ is \emph{below} the majority baseline $1.0000$
and sign criteria (i)/(ii) FAIL; yet $12/12$ training-free features are
BH-significant (strongest \texttt{reach\_b} $\rho\!\approx\!+0.916$,
\texttt{motif\_edges\_kept} $\rho\!\approx\!+0.905$) and magnitude MAE
$0.0840$ beats the constant baseline $0.2310$, so (iii)/(iv) PASS. The
magnitude association holds while sign prediction degenerates for lack of
sign variance; the prediction is partially supported (magnitude), not
confirmed as a sign predictor. Third, a \emph{fully anticipatory} preregistered replication (C1c) on two
independently constructed real MaleCNS subgraphs under different visual-ID
rules, each with $|P_{\mathrm{all}}|=3$ and four distinct dose levels
($p_2\!\neq\!p_{\mathrm{full}}$), reproduces the monotone dose on both
(\texttt{cross\_subgraph\_replicated=true}), removing C1b's remedial-design
and two-dose limitations. It remains a controlled diagnostic task:
\texttt{mechanism\_transfer\_confirmed} stays false,
\texttt{expected\_strong} is still unmet, and no improved visual accuracy is
claimed.
\end{abstract}

\input{tables}
\input{tables_taskaware}
\input{tables_r39}
\input{tables_r46}
\input{tables_r53}
\input{tables_r56}
\input{tables_r59}
\input{tables_r60}

\section{Introduction}
\label{sec:intro}

Connectome-constrained models promise a biological wiring diagram as an
inductive bias, yet reports disagree on whether it helps: fixed-connectivity
learned-gain models report efficient visual computation at matched wiring
cost~\cite{structurealone2026}, while a strict degree-preserving null with
shared initialisation removes much of the topological
advantage~\cite{toposens2026}; motor-domain and lesion-audit results add to
the conflict~\cite{flygm2026,flyhack2026}. These studies are not comparable --
they differ in null family, training horizon, task, and above all in whether
the task is \emph{structure-sensitive at all}. The field lacks a comparison
making the null model itself the experimental variable; we treat the control
generator as a measurable object validated by controlled causal interventions.
\emph{Scope}: an evaluation-methodology study whose real-image probes are
diagnostics, not a benchmark; the XOR label is not a natural visual task and
we claim no improved visual accuracy.

We preregister a \emph{nested ladder} of nulls matching, in order, edge count
(L0), per-node in/out degree (L1), module-pair counts (L2) and the full
directed triad census (L3), with an independent same-generator draw (L4) as a
matched control; every rung carries invariant assertions. The label is an XOR
of two distant sources solvable only through an injected two-hop convergence
motif, so removing the graph (or those edges) removes the only information
path. Ten frozen seeds, paired bootstrap intervals, Benjamini--Hochberg
correction~\cite{benjamini1995} and leave-one-graph-out evaluation complete
the protocol. Three outcomes were admissible \emph{a priori}: a localising
monotone ladder, a preregistered null result, or a diagnostic outcome in
which the ladder itself is inadequate.

We observe outcome (iii): the effect is large but every rung erases it alike.
We then extend the diagnosis to real wiring and the published baseline.
\emph{Scope}: within controlled low-in-degree graphs and a specific real
MaleCNS subgraph regime, path-aware control construction decides the observed
structure effect and the published weak/strict ladder is reproduced on that
real data; \texttt{mechanism\_transfer\_confirmed=false} precludes any causal
explanation of the literature disagreement, and real visual tasks are negative
for topology specificity. Our contributions:

\begin{itemize}
\item A \emph{reusable null-generator standard}: nested rungs with
  machine-checked invariants, a structure-sensitivity gate, corrected
  multi-seed statistics, a task-aware null and motif-dose ladder closing the
  controlled-graph causal loop, and a clean hash-verifiable 9-condition
  $\times$ 10-seed rerun passing all six gates (Table~\ref{tab:c5gates}).
\item A \emph{path-level dose--response on real MaleCNS v1.0 wiring}: a
  preregistered fixed-density motif dose fails
  (\texttt{hypothesis\_supported=false}) because degree-preserving swaps
  create fresh two-hop paths that global invariant assertions miss; a
  preregistered path-blocking repair yields a monotone dose with all six
  rules satisfied (\texttt{hypothesis\_supported\_pathblock=true},
  Table~\ref{tab:dose}).
\item A \emph{fully anticipatory} preregistered replication on two
  independently constructed real MaleCNS subgraphs (\S\ref{sec:c1c}), each with
  $|P_{\mathrm{all}}|=3$, four dose levels and all seven rules satisfied
  (\texttt{cross\_subgraph\_replicated=true}), removing C1b's remedial-design
  and $p_2{=}p_{\mathrm{full}}$ limitations without leaving the controlled
  diagnostic regime.
\item A faithful rerun of the published weak/strict/task-aware control ladder
  reproducing the weak/strict pattern on this real data
  (\texttt{mechanism\_transfer\_confirmed=false}) and a connectome-constrained
  model whose gain comes from having a graph, not the specific topology -- an
  equivalent re-implementation, MNIST not run; we therefore cannot claim to
  explain the causal mechanism of the literature disagreement.
\item Full disclosure of negative results and failures: the \textsc{failed}
  three-seed pre-check, dense chance-level runs, the
  \texttt{hypothesis\_supported=false} v2 verdict, the negative real-image
  boundary, the completed but negative 24/24 independent family
  (\texttt{independent\_family\_validated=false}: sign criteria FAIL,
  magnitude criteria PASS), and an explicit cross-run visual-ID comparability
  warning.
\end{itemize}

\section{Related Work}
\label{sec:related}

\paragraph{Connectome-constrained vision.}
Lappalainen et al.~\cite{lappalainen2024} introduced a cell-type-resolved
connectome-constrained network with an optic-flow task; later work fixed the
graph and learned only synaptic gains or a linear
readout~\cite{structurealone2026}, and whole-brain
message-passing controllers were studied in the motor
domain~\cite{flygm2026}. The public MaleCNS v1.0~\cite{malecns2026} and
FlyWire~\cite{dorkenwald2024} connectomes are separate, non-interchangeable
resources with different neuron and synapse counts.

\paragraph{Null models and network motifs.}
Degree-preserving randomisation and edge swaps separate topology from
lower-order constraints~\cite{maslov2002,milo2002}; the sixteen-class directed
triad census~\cite{wasserman1994} is implemented and cross-checked here
against NetworkX~\cite{hagberg2008}. The closest prior
work~\cite{toposens2026} applies a degree-preserving null with shared
initialisation and reports that early-training advantages largely disappear;
its limitations name longer horizons, broader tasks and higher-order nulls as
future work -- the axes our ladder varies. Ensemble-level
interventions~\cite{flyhack2026} and reconstruction-variation
audits~\cite{reconvar2026} address orthogonal questions.

\paragraph{Evaluation protocol.}
Our task follows synthetic structure-sensitivity tests in graph learning, where
a label is solvable only through a specified relational
path~\cite{kipf2017,xu2019}. We additionally require a leakage probe at chance
and a no-graph ablation that collapses the model to chance, distinguishing
\emph{structural} advantage from mere capacity.

\section{Method}
\label{sec:method}

\begin{figure}[t]
  \centering
  \includegraphics[width=\linewidth]{figures/fig1_framework_gen.png}
  \caption{Overall framework. Inputs (a controlled two-hop XOR graph, a real
    MaleCNS v1.0 visual subgraph, digit images) pass through a nested null
    ladder (L0--L3 invariants, matched control, path-survival audit) into a
    shared-init two-round message-passing backbone with frozen seeds and
    oracle/leakage probes. Preregistered tests (path dose and path-blocking
    repair, published control ladder, connectome visual model, clean rerun,
    24-graph family) yield held-out accuracy with paired bootstrap intervals
    and BH correction. Dashed box: scope conditions.}
  \label{fig:framework}
\end{figure}

Figure~\ref{fig:framework} summarises the design: the null generator, not the
graph family, is the experimental variable, and every condition shares the
same backbone, initialisation and seeds.

\subsection{Controlled graph and data provenance}
\label{sec:graph}

We use a 96-node directed graph with eight modules ($p_{\text{in}}=0.06$,
$p_{\text{out}}=0.004$) and four injected two-hop convergence motifs
$(\text{sources } a,b)\to \text{relay } r\to \text{sinks } S$, with all
$a\!\to\!S$ and $b\!\to\!S$ edges removed (123 edges). Module sizes
$[25,17,12,11,11,7,7,6]$ come from the top-eight \texttt{class} distribution
of the public 14.5\,MB MaleCNS v1.0 annotation file
(CC-BY~4.0~\cite{malecns2026}). \emph{Scope condition}: this is a
\textsc{synthetic\_approximation}. We did not parse the 1.05\,GB full
connectome weight matrix -- an earlier attempt was silently terminated by an
out-of-memory kill -- so no claim here concerns the real MaleCNS visual
subgraph; the only biological input is the module-size prior. All effects
below are properties of the controlled graph and protocol until
\S\ref{sec:realconnectome}.

\subsection{Structure-sensitive task}
\label{sec:task}

Each sample draws $X\in\mathbb{R}^{N\times d}$, $d=16$. Source $a$ carries a
signal in the first $d/2$ channels, source $b$ in the rest; all other channels
and nodes carry Gaussian noise (scale $0.06$). With fixed unit $u,v$,
\begin{equation}
y = \mathbb{1}[x_a\!\cdot\! u > 0]\;\oplus\;\mathbb{1}[x_b\!\cdot\! v > 0].
\end{equation}
The label is not linearly decodable from raw or mean-pooled features; the
disjoint informative channels let a sum-aggregating relay retain both signs.
A two-layer network with in-neighbour normalisation
$W = A^{\top}\mathrm{diag}(A^{\top}\mathbf{1})^{-1}$ and sink readout has
exactly the two-hop receptive field required. \emph{The in-neighbour direction
is essential}: using $A$ instead of $A^{\top}$ makes the relay aggregate its
out-neighbours and the real graph collapses to chance; we found and corrected
this defect in an earlier run (\S\ref{sec:gate}). Three controls are
preregistered: an $A=0$ reservoir, a motif-only upper-bound graph, and a
shuffled node-to-input mapping.

\subsection{Nested null ladder and invariant assertions}
\label{sec:ladder}

Each rung generates independent nulls and logs per-instance invariant
violations: \textbf{L0} same edge count (Erd\H{o}s--R\'enyi); \textbf{L1} plus
exact in/out-degree by double-edge swaps; \textbf{L2} plus exact module-pair
edge counts by within-cell swaps; \textbf{L3} plus the sixteen-class directed
triad census by a Metropolis chain rejecting census-worsening swaps (residual
error logged, never hidden); \textbf{L4} (a matched \emph{control}, not a
rung) an independent draw from the same generator. We generate 16 instances
for L0--L2/L4 and 8 for L3 (600 Metropolis steps); the triad census uses a
vectorised 64-pattern lookup cross-validated against
NetworkX~\cite{hagberg2008}. We report maximum violations per rung. L4 is
excluded from monotonicity tests because it violates every invariant -- a
correction to an earlier protocol that treated it as the strictest rung. The
primary statistic is the two-path count
$\mathrm{conv2}(A)=|\{(i,j),i\ne j:(A^2)_{ij}>0\}|$, equal to 178 for the
real graph.

\subsection{Why global invariants cannot protect a local path}
\label{sec:theory}

The ladder's failure to localise the effect is structural. Let
$P=\{(a,r),(b,r)\}\cup\{(r,s):s\in S\}$ be the $|P|=18$ injected edges. A
double-edge swap touches a protected edge with probability $2/E$, $E=123$;
after $T$ successful swaps,
\begin{equation}
\Pr[e\in A_{\text{null}}]\approx\left(1-\frac{2}{E}\right)^{T}\le e^{-2T/E},
\label{eq:survival}
\end{equation}
below $10^{-7}$ per edge for the $T\!\approx\!2\!-\!5\!\times\!10^{3}$ swaps
of L1--L3. The label needs the \emph{conjunction} of a source-to-relay and a
relay-to-sink edge, so losing any member removes the only path; restricting
swaps to module cells or rejecting census-worsening ones changes \emph{which}
swaps are admissible but never protects $P$. Equation~\eqref{eq:survival}
assumes uniform swaps, so for L2/L3 it is an upper-bound intuition, not an
exact per-rung probability; we therefore measure surviving task-relevant edges
directly (\S\ref{sec:taskaware}). Any sufficiently mixed null drives accuracy
to chance regardless of matched global statistics (consistent with the
$0.0244$ rung spread), and monotonicity across rungs is not well-posed -- rung
order measures global fidelity, while the task depends on a local subgraph
orthogonal to it. A control is sufficient only if it preserves every structure
the hypothesis does not ablate.

\subsection{Training, statistics and held-out prediction}
\label{sec:protocol}

All conditions share the architecture (hidden width 64, two propagation
rounds, sink mean-pooling, one linear head), initialisation,
$n_{\text{train}}=n_{\text{test}}=512$, 120 epochs, Adam at
$3\!\times\!10^{-3}$/$6\!\times\!10^{-3}$ and ten frozen seeds
$\{11,17,23,31,41,53,67,79,97,103\}$; test stimuli are fixed and never used
for tuning. We report paired seed-level differences
$\Delta_c=\mathrm{acc}(\text{real})-\mathrm{acc}(c)$ with 5000-resample
percentile bootstrap 95\% intervals, a two-sided bootstrap $p$ and
Benjamini--Hochberg~\cite{benjamini1995} $q$ across the eight-condition
family. For held-out prediction we generate five independent graphs
(seeds 211/307/401/503/601), compute training-free features (two-path counts,
two-hop reachability, module density and differences against nulls), and a
leave-one-graph-out loop reselects the best-correlated feature on four
training graphs to predict the held-out delta, reporting sign agreement and
magnitude MAE against a constant baseline.

\section{Experiments and Results}
\label{sec:exp}

\subsection{Setup and hardware}
\label{sec:setup}

All experiments ran remotely on one NVIDIA RTX~4090 (46\,068\,MiB, driver
580.126.09), Python 3.11.18, PyTorch 2.11.0+cu130, CUDA 13.0, with a verified
CUDA smoke test; total wall time 151\,s. Configurations, seeds, per-seed
accuracies, invariant counts and logs are persisted with the run.

\subsection{Structure-sensitivity gate: an honest failure}
\label{sec:gate}

A three-seed pre-check recorded \textsc{failed}: its auxiliary oracle clause
($>0.80$) scored $0.7891$, while all other clauses passed (mean real
$0.8092$, paired CI lower bound $0.2988$, $p=0.0002$, clean motif $0.8470$,
leakage probes $0.4902$/$0.5000$). We do not relabel this verdict. Retraining
the probe for 600 epochs still gives $0.7930$ (32 concatenated channels, only
two carry the label; a clean two-channel oracle reaches $0.9232$ against a
noise-implied ceiling of $0.9395$). The failing clause is a probe defect, not
a precondition failure; the ten-seed main matrix below -- with its $A=0$
control, upper-bound path and falsifying ablation -- is the primary evidence.

\subsection{Main matrix}
\label{sec:main}

\begin{figure}[t]
  \centering
  \includegraphics[width=0.94\linewidth]{figures/fig2_main_results.pdf}
  \caption{Ten-seed held-out accuracy. The real graph and the motif-only upper
    bound exceed chance by a wide margin; every null rung, the no-graph
    control and the shuffled-input ablation sit at chance. Error bars are
    seed standard deviations.}
  \label{fig:main}
\end{figure}

Table~\ref{tab:main} and Fig.~\ref{fig:main} summarise the campaign. The real
graph reaches $0.8152$ (seed variance $2.1\!\times\!10^{-4}$); the no-graph
control, all four rungs and the shuffled ablation lie between $0.4795$ and
$0.5133$. Every real-minus-condition difference is $0.30$--$0.34$ with 95\%
intervals excluding zero and BH $q=0.0002$. The motif-only graph slightly
exceeds the full real graph ($0.8348$, $\Delta=-0.0195$): the injected motif
is the information path, background edges add noise. Shuffling node-to-input
mapping collapses accuracy to $0.5133$, ruling out a readout artefact. The
decisive observation (last column of Table~\ref{tab:main}) is that the four
rungs differ by at most $0.0244$ in their real-minus-rung delta, an order of
magnitude below the effect: \textbf{the ladder has no localising power}.
Matching degree, module blocks and triad census removes the advantage as much
as a plain edge-count-matched random graph.

\subsection{Ladder statistics and invariant verification}
\label{sec:invariants}

Table~\ref{tab:invariants} reports maximum invariant violations: L1 preserves
degree exactly, L2 every module cell, and L3 reduces the triad-census $L_1$
deviation to 20 from 1826 at L2; L4 violates every invariant (28 edges, 93
degree-mismatched nodes, 39 module cells), confirming its reclassification.

\begin{figure}[t]
  \centering
  \includegraphics[width=0.98\linewidth]{figures/fig3_ladder_analysis.pdf}
  \caption{(a) Absolute mean difference of the two-path statistic from the
    real graph is non-monotone across L0--L3 (one inversion). (b) Paired
    accuracy differences with 95\% bootstrap intervals; all null conditions
    are statistically indistinguishable.}
  \label{fig:ladder}
\end{figure}

The two-path statistic (Fig.~\ref{fig:ladder}a) has absolute mean deviation
$24.875$, $0.125$, $5.938$, $0.500$ for L0--L3 -- non-monotone with one
inversion -- and paired accuracy differences (Fig.~\ref{fig:ladder}b) show all
nulls statistically indistinguishable. Degree preservation alone already
reproduces the \emph{global} two-path count while accuracy on L1 is at chance:
matching a graph statistic is not preserving the computational path.

\subsection{Robustness}
\label{sec:robust}

Table~\ref{tab:robust} reports $\mathrm{acc}(\text{real})-\mathrm{acc}(\text{L1})$
under two noise levels and two propagation depths. The effect persists in all
four cells ($0.1406$--$0.3652$), degrading gracefully with noise; we retain
all four runs, including the weakest.

\subsection{Powered held-out prediction}
\label{sec:heldout}

An initial five-graph study produced only positive deltas, making sign
prediction unidentifiable; we therefore preregistered a powered follow-up with
\textbf{twenty} held-out graphs (ten intact, ten with 75\% of motif edges
removed). Observed deltas now span $-0.046$ to $0.372$ (12 positive, 8
non-positive). Twelve training-free features were fixed in advance and a
leave-one-graph-out loop reselects the best-correlated feature inside each
fold; the script reports its own SHA-256, verified against the pulled file.
Eight of twelve features remain significant after correction (two-path count,
differential reachability, edge/module density, surviving motif-edge count and
fraction, both source reachabilities; Table~\ref{tab:heldout}). LOGO magnitude
prediction reaches MAE $0.1046$ vs a constant-mean baseline $0.1635$ (36\%),
but \textbf{sign agreement is $0.60$ -- exactly the majority-class baseline}:
the sign remains unpredictable. We claim only a corrected association between
training-free statistics and magnitude, not a causal or sign-level predictor.
Because half the graphs were deliberately motif-degraded, several significant
features partly encode the manipulation itself; establishing predictive power
beyond it requires a separate, non-manipulated family, which we supply as C4
(\S\ref{sec:c4family}): there the magnitude association survives on 24
non-manipulated graphs, while sign prediction again degenerates.

\subsection{Task-aware null and motif dose--response}
\label{sec:taskaware}

\S\ref{sec:theory} predicts a control is sufficient only if it preserves the
structures the hypothesis does not ablate. We test this constructively in a
companion 10-seed run: a \emph{task-aware null} protects all 18 motif edges
and randomises only the background under degree/module constraints; a
\emph{strong null} also removes all motif edges; a dose ladder deletes
$0,4,9,14,18$ of the 18 edges from the real graph. All conditions share the
frozen seeds.

\begin{figure}[t]
  \centering
  \includegraphics[width=0.94\linewidth]{figures/fig4_dose_response.pdf}
  \caption{Motif dose--response. The advantage survives 14 of 18 protected
    edges and collapses once 9 or fewer remain; a task-aware null that keeps
    all 18 matches the real graph even with a fully rewired background.}
  \label{fig:dose}
\end{figure}

Table~\ref{tab:taskaware} and Fig.~\ref{fig:dose}: the task-aware null reaches
$0.8346$, indistinguishable from real ($0.8371$; $\Delta=0.0025$, CI
$[-0.0025,0.0084]$, $q=0.40$) despite a fully rewired background; the strong
null falls to $0.4920$ and motif removal to $0.5117$ ($q=0.00032$). The dose
ladder is sharply thresholded ($0.8371$, $0.8316$, $0.4988$, $0.5104$,
$0.5117$ for $18,14,9,4,0$ surviving edges), and a background-only null with
14 edges intact ($0.8260$) still shows the advantage -- insufficient mixing
yields a control that does not ablate the hypothesis. The advantage is carried
by a small identifiable subgraph: \emph{how many} task-relevant edges a
control preserves, not its rung, determines the effect;
Eq.~\eqref{eq:survival} explains why ordinary generators destroy them. This
companion run also cross-validates \S\ref{sec:main} (real $\approx\!0.84$ vs
no-graph/strong-null $\approx\!0.49$--$0.51$, paired CIs excluding zero,
$q\!\approx\!0.0003$) without erasing the execution defect disclosed in
\S\ref{sec:discussion}.

\subsection{External validation on a real MaleCNS visual subgraph}
\label{sec:realconnectome}

To test transfer to biological wiring, we streamed the 1.05\,GB MaleCNS v1.0
weight table (CC-BY~4.0~\cite{malecns2026}), selected 9{,}766 visual IDs and
formed a 768-node induced subgraph (48{,}713 edges after ten shortcut removals;
48{,}723 before); the matrix was never densified. A genuine two-hop convergence
motif needed no injection. Over six frozen seeds, real / $A\!=\!0$ /
degree-preserving / task-aware / motif-removed all stay at chance
(Table~\ref{tab:real}), every paired interval containing zero: \textbf{the
structure-sensitivity precondition of \S\ref{sec:task} fails on this dense
subgraph}. With mean in-degree $\approx\!63$, propagation likely dilutes the
signals; no oracle control was run here, so this stays a hypothesis until
\S\ref{sec:density}. We report this boundary in full and scope the controlled
conclusion to low-in-degree, path-identifiable subgraphs.

A preregistered follow-up sparsified the \emph{same} subgraph to top-4
in-edges ($48{,}723\to3{,}072$ edges, in-degree $63.4\to4.0$; $3{,}062$ after
shortcut removal). \textbf{Five settings changed together} -- top-4
sparsification, readout mean$\to$sum, width $64\!\to\!128$, lr
$0.006\!\to\!0.008$, epochs $120\!\to\!300$ -- with seeds/graph/motif fixed, so
this cannot isolate density, only establish a solvable regime. The precondition
is now met (Table~\ref{tab:real}, ``Sparse''): real $0.6423$ vs $A\!=\!0$
$0.5283$ ($\Delta=0.1139$, CI $[0.0485,0.1680]$), motif removal $0.5160$
($\Delta=0.1263$, CI $[0.0736,0.1608]$). Yet the degree-preserving null
\emph{retains} the advantage ($0.6162$; $\Delta=0.0260$, CI
$[-0.0348,0.0921]$) and task-aware gives $0.5850$ ($\Delta=0.0573$, CI
$[0.0059,0.1029]$). The preregistered motif-survival-threshold hypothesis is
\textbf{not} confirmed (\texttt{hypothesis\_supported=false}); we make no
replication claim -- whether a null removes the effect depends on target-graph
density, not the rung label, which supports measuring structure sensitivity
and survival counts \emph{on the target graph}.

\subsection{Isolating density on real wiring}
\label{sec:density}

The five-setting follow-up cannot attribute the change, so we ran a
\emph{single-variable} sweep on a freshly streamed subgraph (768 nodes,
44{,}815 edges, in-degree 58.35; 29{,}442 visual IDs), varying only $k$, the
number of retained in-edges; readout, epochs (120), width (64), lr (0.006),
six seeds and motif are fixed. An oracle MLP on both source features reaches
$1.0000$: dense failure is a propagation, not a label, problem.
Table~\ref{tab:density} shows a sharp threshold: real stays at chance for
in-degree 58.4/49.0/15.9 ($0.5049$/$0.4886$/$0.5075$), rises to $0.6227$ at
8.0 and reaches $0.9863$/$0.9987$ at 4.0/2.0; $A\!=\!0$ and degree-null stay
near $0.50$, task-aware recovers in parallel ($0.9027$/$0.9971$ at
$k\!=\!4/2$). \textbf{Density alone determines whether the gate passes}
(\texttt{density\_isolated=true}); it does \emph{not} establish motif survival
as causal -- at $k\!=\!4,2$ the protected edge $23\!\to\!63$ is absent yet real
is near ceiling. Visual-ID matching differs across runs (9{,}766 / 29{,}442 /
29{,}564 IDs), so absolute counts are not comparable with
\S\ref{sec:realconnectome}; within the sweep all $k$ share one graph and motif.

\subsection{A real visual task: a negative boundary}
\label{sec:visualtask}

To test whether the connectome prior is a competitive visual bias, we ran real
image classification (\texttt{sklearn} digits, $8\times8$, 10 classes, fixed
1078/359/360 split; plus an MNIST add-on) through a 768-node MaleCNS subgraph,
comparing real, degree-preserving null, $A\!=\!0$, task-aware null and a
non-graph MLP over six seeds. Real beats no graph ($+0.0995$, CI
$[0.0556,0.1648]$), but the degree-preserving null matches it exactly
($-0.0005$, CI $[-0.0412,0.0440]$) and a plain MLP dominates every graph
model ($0.9722$ vs $0.7185$; MNIST $0.8879$ vs $0.7649$). \textbf{We report
this as a negative result}: on these tasks the wiring prior is not a
competitive bias, and the gain over no graph is already explained by the
degree-preserving control. On the MNIST add-on, however, real \emph{does} beat
the degree-preserving null by a small significant margin ($+0.0148$, CI
$[0.0100,0.0192]$), opposite in sign to the digits result; we report both --
null behaviour is dataset-dependent, which further supports running control
comparisons on the target task. The digits column of Table~\ref{tab:visual}
also lists the task-aware null ($\Delta=-0.0148$, CI $[-0.0278,-0.0009]$),
itself significant and negative; we state it here rather than leaving it only
in the table.

\subsection{Faithful rerun of the published control ladder on real wiring}
\label{sec:ladderrerun}

A diagnosis shown only on a controlled surrogate could be dismissed as its
artefact. We therefore reran the three-segment control ladder
of~\cite{toposens2026} (arXiv:2604.04033) on a real MaleCNS v1.0 visual
subgraph (768 nodes / 41{,}322 edges / 105{,}271 visual IDs; mean in-degree
53.80), six frozen seeds, with one shared GNN whose parameters are identical
across conditions (Table~\ref{tab:r53ladder}). The weak segment, resumed from a
real-graph checkpoint, reproduces the apparent early advantage: null-minus-real
loss at step~5 is $+0.1454$, CI $[+0.1030,+0.1892]$. The strict
shared-initialisation segment abolishes it ($\approx 0$, CI
$[-0.0003,+0.0003]$). Our task-aware extension, protecting the two-hop paths
from the input neurons, does \emph{not} restore it: the preregistered
difference-in-differences is $-0.0477$, CI $[-0.1056,+0.0120]$. We therefore
report \texttt{mechanism\_transfer\_confirmed=false}. The $k{=}4$ robustness
variant (within-task-aware gap $+0.0870$, CI $[+0.0394,+0.1366]$) does not
gate the decision and is not used to rescue it. This is an equivalent
re-implementation, not the official flyvis stack.

\subsection{A connectome-constrained visual model on real images}
\label{sec:constrained}

To test whether the wiring prior helps on real vision data, we trained a
fixed-wiring model with learnable per-synapse scalar gains on the same
subgraph, reading out a linear classifier on sink activities
(Table~\ref{tab:r53constrained}). On ten-class real $8\times8$ digit images,
real wiring reaches $0.9431$ while the no-graph control stays at $0.1028$
(paired gap $+0.8403$, CI $[+0.8370,+0.8444]$), satisfying criterion~(iii).
The degree-preserving null, however, is slightly \emph{higher} ($0.9509$;
gap $-0.0079$, CI $[-0.0120,-0.0042]$) and the task-aware null matches real
exactly: the benefit comes from having a sparse directed graph at all, not
from the specific biological topology -- an honest negative for any
topology-specific claim. This is an \emph{equivalent re-implementation}, not
the official flyvis code; the optional MNIST extension was preregistered but
not run, so no MNIST claim is made.

\subsection{Fixed-density path-survival dose on real wiring}
\label{sec:pathblock}

The controlled mechanism predicts that deleting every two-hop computation
path -- holding edge count, degree sequence and module-pair counts fixed --
must collapse accuracy to chance. We preregistered this fixed-density dose on
real MaleCNS wiring (Table~\ref{tab:dose}). The first run (\emph{C1}) returned
\texttt{hypothesis\_supported=false}: four of six rules failed and the real
dose was non-monotone (full $0.9772$, $s_2/s_1$ $0.5051$, $s_0$ $0.9745$).
Code-level diagnosis: after deleting a nominal motif edge, the
module-pair-preserving density swaps were not forbidden from creating a
\emph{new} common-relay triple, so fresh two-hop paths appeared while every
global invariant passed (rule~vi held). A second run (\emph{C1b}) was
preregistered to change only this: reject any swap creating a new common-relay
triple; seeds, hyperparameters, conditions, statistics and the six-rule
structure are unchanged, and the original negative verdict is retained as a
read-only contrast. C1b returns
\texttt{hypothesis\_supported\_pathblock=true} with all six rules true: real
rises $0.5055 \to 0.9641 \to 0.9772$ as $0 \to 1 \to 2$ paths survive, while
degree-null and $A{=}0$ stay at chance and leakage probes $\approx 0.505$.
Because $|P_{\mathrm{all}}|{=}2$ on this subgraph, $p_2$ and
$p_{\mathrm{full}}$ are identical -- reported, not hidden. These paired runs
establish that \textbf{global invariant assertions are insufficient to exclude
a path-level confound}; control construction must be path-aware.

\subsection{Independent real-wiring graph family (C4)}
\label{sec:c4family}

The powered held-out study (\S\ref{sec:heldout}) left open whether its
significant features predict beyond the intact/degraded manipulation. C4 is
the \emph{first real test} of that second falsifiable prediction, on
\textbf{24/24 completed} non-manipulated real-wiring graphs: four graph seeds
$\times$ six top-$k$ sparsities, 12 training-free features, LOGO-CV selection
inside training folds. Artifacts are on disk
(\texttt{results/r51\_independent\_family/}) and pass
\texttt{verify\_c4\_consistency.py} (frozen config SHA
\texttt{aa9c19c6\ldots}). The verdict
\texttt{independent\_family\_validated=false} is mixed, partly negative
(Tables~\ref{tab:c4family}--\ref{tab:c4bh}).

Sign side: all 24 deltas are non-negative (eight high-$k$ graphs
$k\!\in\!\{8,12\}$ near $+0.50$, the rest $\approx0$, one $+0.0037$), so no
sign variance exists. LOGO sign agreement $0.5000$ is \emph{below} the
majority baseline $1.0000$; criteria (i) \texttt{i\_signs\_not\_all\_same}
and (ii) \texttt{ii\_sign\_agreement\_margin} FAIL. Magnitude side: all
$12/12$ features are BH-significant ($q<0.05$; \texttt{reach\_b}
$\rho\!\approx\!+0.916$, \texttt{motif\_edges\_kept} $\rho\!\approx\!+0.905$),
and LOGO MAE $0.0840$ beats the constant baseline $0.2310$, so (iii)/(iv)
PASS. The prediction is \emph{partially supported} (magnitude only): not
confirmed as a sign-level predictor, with the preregistered
\texttt{independent\_family\_validated} verdict false.

\subsection{Fully anticipatory cross-subgraph replication (C1c)}
\label{sec:c1c}

C1b was preregistered \emph{after} inspecting C1's in-flight readings (a
remedial design), and its $|P_{\mathrm{all}}|=2$ subgraph gave
$p_2=p_{\mathrm{full}}$, leaving only two accuracy levels. C1c removes both: a
complete protocol was frozen \emph{before any training} on two independently
constructed real MaleCNS subgraphs under different visual-ID rules, each with
$|P_{\mathrm{all}}|\geq3$ and $\geq3$ distinct dose levels, retaining the
path-blocking swap, 10 paired seeds, five conditions and the six-rule
structure. Both reproduce the monotone dose (Table~\ref{tab:c1cdose}): real
$0.5017\!\to\!0.6889\!\to\!0.8955\!\to\!0.9432$ (broad) and
$0.5017\!\to\!0.6742\!\to\!0.8877\!\to\!0.9407$ (narrow), while $A{=}0$ stays
at $0.5017$; $\Delta(\mathrm{full}-p_0)=+0.4415$ [0.426, 0.454]
($p_{\mathrm{BH}}=0.0025$) and $+0.4390$ [0.425, 0.453]
($p_{\mathrm{BH}}=0.0036$). All seven rules pass on both
(\texttt{cross\_subgraph\_replicated=true}).

Limitations are reported, not hidden: a \texttt{KeyError} crashed the first
run, and its fix changed one execution script's SHA while the verified
preregistered config, thresholds, seeds and graph-construction seed did not
(failed log retained); rule (iii) is a \emph{grand mean pooled over doses}, so
individual null cells exceed $0.60$ (task-aware null $0.9432$ at full dose by
design, one degree-null cell $0.6742$) and nulls are not everywhere at chance;
both subgraphs share the MaleCNS v1.0 release and a controlled XOR label, not
a natural visual task, so they are independently \emph{constructed}, not
independent datasets. C1c strengthens the path-aware-control finding
\emph{within this regime} without establishing mechanism transfer
(\texttt{mechanism\_transfer\_confirmed} remains false), reversing earlier
negative results, or meeting \texttt{expected\_strong}.

\section{Discussion}
\label{sec:discussion}

\paragraph{Scope of the claim.}
Evidence is scoped to controlled low-in-degree graphs and two real MaleCNS
subgraph regimes within one release: path-aware control construction decides
the observed structure effect and the published weak/strict ladder is
reproduced there. \texttt{mechanism\_transfer\_confirmed=false} precludes any
causal explanation of the conflicting literature; topology specificity is
negative on real visual tasks (\S\ref{sec:visualtask}, \S\ref{sec:constrained})
and C4's sign prediction fails (\S\ref{sec:c4family}). We make no cross-study
causal claim.

\paragraph{Distinct preregistered tests with different outcomes.}
Several decisive results must not be merged or quoted selectively. The faithful
ladder rerun (\S\ref{sec:ladderrerun}) returns
\texttt{mechanism\_transfer\_confirmed=false} (its task-aware segment does not
recover the advantage), while the fixed-density dose
(\S\ref{sec:pathblock}) returns
\texttt{hypothesis\_supported\_pathblock=true} only after an adaptive
remedial re-preregistration (C1b), with the original C1 negative verdict
retained read-only. These are \emph{different} tests reported both, quoting
neither selectively. The independent-family test
(\S\ref{sec:c4family}) is a third, different verdict
(\texttt{independent\_family\_validated=false}), never merged into either.
C1c (\S\ref{sec:c1c}) is a fourth, \emph{fully anticipatory} replication on
two real subgraphs (\texttt{cross\_subgraph\_replicated=true}) removing C1b's
remedial-design and two-dose limitations; its positive verdict does not
retroactively confirm C1, reverse C4's negative sign verdict, or rescue
\texttt{mechanism\_transfer\_confirmed=false}. The conjunctive
\texttt{expected\_strong} target is unmet; no favourable subset is presented
as the study's outcome.

\paragraph{Relation to the strongest published control ladder.}
The closest prior work~\cite{toposens2026} reports an early-training advantage
that disappears under shared initialisation and a degree-preserving null. Our
protocol \emph{begins} at that strict endpoint and adds the axes its
limitations name -- higher-order rungs, full convergence, corrected statistics,
a predictive statistic -- plus a record of preserved task-relevant edges. We
rerun that ladder on real wiring (\S\ref{sec:ladderrerun}), reproducing the
weak advantage and its strict abolition; ours is an equivalent
re-implementation, not the official flyvis stack. On the controlled graph a
task-aware null restores the advantage (\S\ref{sec:taskaware}) -- path, not
degree, preservation matters -- yet on real wiring it does not, bounding how
far that claim transfers.

\paragraph{Using the standard in practice.}
Before interpreting a ``topology helps'' claim: (1) prove structure
sensitivity with a no-graph control and leakage probe; (2) record which
task-relevant structures each control preserves (Eq.~\eqref{eq:survival}
shows unprotected edges vanish under mixing, so this count, not the rung
label, is informative); (3) report survival counts beside the delta; (4)
require paired multi-seed intervals with correction. A control is
\emph{sufficient} only if it preserves every structure the hypothesis does not
ablate and ablates exactly the one it does.

\paragraph{Limitations.}
\emph{Data}: main results use a 96-node controlled graph; real 1.05\,GB
MaleCNS runs are boundaries, not replications. \emph{Baselines}: the published
ladder and connectome-constrained model are equivalent re-implementations,
\emph{not} official flyvis; signed/typed edges unmodelled, MNIST not run.
\emph{Task}: the XOR label is not a natural visual task; real-image probes are
negative for topology specificity (a non-graph MLP dominates; gain is from
having a graph, not its topology), and we claim no improved visual accuracy.
\emph{Statistics}: sign prediction equals the majority baseline, so only
magnitude association is supported. \emph{Protocol}: the three-seed pre-check remains \textsc{failed} with a
documented probe defect; an early aggregation block crashed after flushing
all data (leave-one-graph-out values were recomputed from persisted rows),
and its executed script does not match the recorded checksum, whereas all
later runs are hash-verifiable. \emph{C4}: the independent 24-graph family
is complete (24/24) and its frozen configuration and artifacts are retained;
\texttt{independent\_family\_validated=false} (sign criteria FAIL, magnitude
criteria PASS), with both halves reported in full.
\emph{Comparability}: real-subgraph runs use different visual-ID matchings
(this batch 105{,}271 IDs / 12.49M edges / 41{,}322 edges; earlier 9{,}766 /
29{,}442 / 29{,}564 IDs, 48{,}713 / 48{,}723 / 44{,}815 edges) -- absolute
counts are not poolable; only within-batch paired comparisons are valid. All
runs, including negative ones, are retained; this is not a vision benchmark.

\paragraph{Reproducibility.}
Frozen seeds, hyperparameters, executed scripts, raw and completed metrics,
per-seed logs, invariant counts and the repair script are preserved; every
table and figure is generated directly from the persisted metric files.
Real-subgraph runs have verified executed-script hashes (dense
\texttt{090d9535\ldots}, sparse \texttt{026bf88e\ldots}, sweep
\texttt{360a3808\ldots}, visual \texttt{28f5744d\ldots}). C5 is a clean
hash-verifiable 9-condition $\times$ 10-seed rerun with all six gates passing;
C1b's frozen configuration (\texttt{8ee6605e\ldots}) retains the original C1
negative verdict as a read-only contrast; C4's configuration
(\texttt{aa9c19c6\ldots}) is verified on all 24 graphs; C1c's configuration
(\texttt{996ab547\ldots}) is independently recomputed. All reported metrics
were cross-checked field by field against the persisted run records, and the
failed three-seed pre-check and transfer-repair artefacts are retained.

\paragraph{Ethics and broader impact.}
Purely computational: no human subjects, animal work or deployed system; the
only biological input is a public CC-BY~4.0 annotation file (credited). An
ill-posed null yielding a confident but wrong structural attribution motivates
reporting survival counts.

\section{Conclusion}
\label{sec:conclusion}

Four distinct preregistered tests are reported together, never merged or
quoted selectively. (1) On controlled low-in-degree graphs a global-statistic
ladder cannot localise an effect a task-aware null recovers; a clean
hash-verifiable rerun passes all six gates. (2) On real MaleCNS wiring a
fixed-density dose failed because density-compensating swaps created fresh
two-hop paths that global invariants missed; an adaptive remedial
path-blocking re-preregistration yielded a monotone dose, and a fully
anticipatory replication on two independently constructed real subgraphs
reproduces the four-level dose on both
(\texttt{cross\_subgraph\_replicated=true}), removing C1b's remedial-design
and $p_2{=}p_{\mathrm{full}}$ limitations. (3) A faithful rerun of the
published ladder reproduced the weak/strict pattern but not its task-aware
recovery (\texttt{mechanism\_transfer\_confirmed=false}); a
connectome-constrained model gained from having a graph, not its topology.
(4) C4 (24/24) returns \texttt{independent\_family\_validated=false}:
magnitude association holds, sign prediction degenerates. The conjunctive
\texttt{expected\_strong} target is unmet; all negative boundaries and
failures are retained. We claim no improved visual accuracy: findings are
scoped to controlled low-in-degree graphs and two real MaleCNS subgraph
regimes within one release, both carrying a controlled XOR label, not a
natural visual task; real visual tasks are negative for topology specificity,
and \texttt{mechanism\_transfer\_confirmed=false} precludes any causal claim.

\clearpage

\section*{LLM Usage Statement}
Large language models were used as implementation and writing assistants.
Claude Code (Anthropic) was used as the agentic coding assistant to draft and
debug experiment scripts and to generate LaTeX tables and figures from
persisted metric files; WorkBuddy with Hunyuan- and Claude-family models was
used for orchestration and literature verification; and Kimi K3 drove the
remote Novix experiment agents that executed the scripts on the stated GPU
hardware and returned the logged artifacts. A separate isolated
general-purpose LLM agent produced the independent review panel used for
internal calibration and the final full-mode acceptance review; these are
machine-generated reviewer simulations, not human peer review. Every
scientific claim, numerical result and bibliographic entry was verified
against executed code, pulled artifacts and public sources; no result or
citation was accepted from a language model without such verification. All
experimental data were generated by the executed scripts on the stated
hardware. The authors take full responsibility for the content.

{
\footnotesize
\bibliographystyle{ieeenat_fullname}
\bibliography{main}
}

\end{document}
